%% P20 --- Optymalizacja stochastyczna i różniczkowanie losowania
%%
%% Zalążek. Poniższy opis jest kontraktem tego programu: co ma uzasadnić i co
%% ma dać czytelnikowi. Napisanie programu oznacza usunięcie bloku
%% \programstub{} i zastąpienie go programem. Jeśli po przerobieniu materiału
%% opis okaże się błędny, popraw go i odnotuj sprzeczność w CLAUDE.md w sekcji
%% *Resolved questions*.

\program{Optymalizacja stochastyczna i różniczkowanie losowania}
\label{prog:P20}

\programstub{%
\textit{[Opis do przetłumaczenia --- patrz wydanie angielskie.]}

Argument: a minibatch gradient is an unbiased estimator with variance
proportional to \code{1/B}, and the noise is a property of the algorithm
rather than a defect in it. Payoff: why the loss curve is noisy and what a
moving average of it hides; the linear scaling rule for batch size and
learning rate, presented as widely repeated folklore with a limited
empirical basis, not as a law; gradient clipping as a bound on the update
rather than on the gradient; and the two ways to get a gradient through a
sampling step --- the score-function estimator (REINFORCE, unbiased and high
variance) and the reparameterisation trick (low variance, requires a
differentiable path) --- which is the fork that separates policy-gradient
RLHF from a VAE.

\medskip
\textbf{Planowana długość:} 50 ramek.
}
